\subsection{Relation to Midge-Swarm and Collective-Motion Studies}

The interpretation is consistent with previous midge-swarm studies that treat
these systems as weakly ordered nonequilibrium collectives with spatial
structure, stochastic components, and effective interactions
\cite{feng2023_sampling,reynolds2017_velocity_gravity,reynolds2018_langevin,reynolds2021_intrinsic,reynolds2024_spatial,sinhuber2021_eos}.
Macroscopic and statistical-physics descriptions make compact state variables
natural explanatory candidates, while Langevin-type studies caution against a
purely deterministic trajectory interpretation. The present analysis adds a
local residual diagnostic: for T1, the tested compact-state representation did
not provide a stable reduction.

The non-uniform spatial sampling result is a useful point of comparison
\cite{feng2023_sampling}. That work made a hidden layer of swarm organization
measurable by showing that individuals do not sample the swarm volume
uniformly. The present analysis makes a different hidden layer measurable by
asking what local motion remains after ordinary local affine geometry has been
subtracted. In both cases, the contribution is to separate structure that can
be obscured by raw trajectories or whole-swarm averages.

The same diagnostic problem may arise in other collective-motion systems, such
as fish schools, bird flocks, cell collectives, or artificial groups. Raw
velocity differences in such systems can reflect whole-group motion, local
deformation, or interaction-driven residuals. Local affine subtraction may
therefore provide a conservative baseline for future comparative studies,
although the present data do not show that the same T1 observable appears
outside midge swarms.
